% ═════════════════════════════════════════════════════════════════════════════
%  chapitres/chapitre1.tex — DevSecOps et Pratiques de Sécurité
% ═════════════════════════════════════════════════════════════════════════════
\chapter{DevSecOps et Pratiques de Sécurité}

\begin{tcolorbox}[
  colback=gray!8, colframe=darkgray,
  arc=4pt, boxrule=0.6pt
]
  Ce chapitre introduit le paradigme DevSecOps et ses fondements conceptuels,
  avant de dresser un panorama structuré des techniques de test de sécurité
  applicables aux environnements modernes. Une attention particulière est portée
  au DAST, dont les enjeux dans les architectures microservices constituent
  la problématique centrale de ce projet.
\end{tcolorbox}

% ─────────────────────────────────────────────────────────────────────────────
\section{Introduction au DevSecOps}
% ─────────────────────────────────────────────────────────────────────────────

\subsection{Du DevOps au DevSecOps : genèse et évolution}

Le mouvement \textbf{DevOps}, apparu au tournant des années 2010, visait à
supprimer le cloisonnement traditionnel entre les équipes de développement
(\textit{Development}) et d'exploitation (\textit{Operations}). En automatisant
les pipelines d'intégration et de déploiement continus (CI/CD), il a permis
d'accélérer considérablement la livraison logicielle tout en améliorant la
stabilité des systèmes.

Cependant, cette accélération a mis en évidence une faille structurelle : la
sécurité restait souvent une préoccupation tardive, traitée en phase finale
par des équipes spécialisées travaillant en silo. Ce modèle, désigné par
l'expression \textit{``Security as a Gatekeeper''}, se révèle incompatible
avec des rythmes de déploiement quotidiens, car il crée des goulots d'étranglement
et laisse des vulnérabilités non détectées atteindre la production.

Le concept de \textbf{DevSecOps} --- contraction de \textit{Development},
\textit{Security} et \textit{Operations} --- émerge comme réponse à cette
inadéquation. Il s'agit d'une culture, d'un ensemble de pratiques et d'outils
visant à intégrer la sécurité à \textit{chaque étape} du cycle de vie logiciel,
en partageant la responsabilité de la sécurité entre toutes les équipes.

\begin{definitionbox}{Définition --- DevSecOps}
  Le DevSecOps est une approche organisationnelle et technique qui intègre la
  sécurité comme une composante intrinsèque et continue du cycle de développement
  et d'exploitation logicielle, par opposition à un contrôle ponctuel en fin de
  cycle. Il repose sur le principe du \textbf{\textit{Shift Left Security}} :
  déplacer les activités de sécurité le plus tôt possible dans le processus
  de développement.
\end{definitionbox}

\subsection{Les principes fondateurs du DevSecOps}

Le DevSecOps repose sur quatre piliers interdépendants :

\begin{enumerate}
  \item \textbf{Automatisation de la sécurité.} Les contrôles de sécurité
    (analyse statique du code, scan de dépendances, tests dynamiques) sont
    intégrés directement dans les pipelines CI/CD et s'exécutent à chaque
    commit, sans intervention manuelle systématique.

  \item \textbf{Collaboration et responsabilité partagée.} Les développeurs,
    les ingénieurs DevOps et les équipes de sécurité (SecOps) travaillent
    conjointement dès la phase de conception. La sécurité n'est plus la
    propriété exclusive d'une équipe spécialisée.

  \item \textbf{Observabilité et retour d'information rapide.} Des mécanismes
    de monitoring, de logging et d'alerting permettent de détecter et de
    corriger rapidement les incidents de sécurité en production.

  \item \textbf{Infrastructure as Code (IaC) sécurisée.} La configuration des
    infrastructures est déclarée sous forme de code versionné, soumis aux mêmes
    revues et analyses de sécurité que le code applicatif.
\end{enumerate}

\subsection{Le Shift Left Security dans les pipelines CI/CD}

La notion de \textit{Shift Left} désigne le déplacement des activités de
vérification vers la gauche du cycle de développement, c'est-à-dire vers les
phases les plus précoces. Cette approche repose sur un constat économique bien
établi : le coût de correction d'une vulnérabilité croît exponentiellement à
mesure que l'on avance dans le cycle de développement.

\begin{figure}[H]
  \centering
  \fbox{\parbox{0.85\textwidth}{\centering\vspace{2cm}
    \textit{[Figure : Courbe du coût de correction selon la phase de détection]}
    \vspace{2cm}}}
  \caption{Le principe du \textit{Shift Left Security} : coût de correction
    en fonction de la phase de détection}
  \label{fig:shift_left}
\end{figure}

Dans une pipeline DevSecOps typique, les contrôles de sécurité s'enchaînent
comme suit :

\begin{itemize}
  \item \textbf{Phase de codage :} analyse statique (SAST), linting de sécurité,
    revue de code assistée.
  \item \textbf{Phase de build :} scan des dépendances (SCA --- \textit{Software
    Composition Analysis}), vérification des images de conteneurs.
  \item \textbf{Phase de test :} tests dynamiques (DAST), tests d'intégration
    de sécurité, IAST.
  \item \textbf{Phase de déploiement :} validation des politiques IaC,
    scan de conformité.
  \item \textbf{Phase d'exploitation :} RASP, monitoring continu, réponse aux
    incidents.
\end{itemize}

% ─────────────────────────────────────────────────────────────────────────────
\section{Les types de tests de sécurité (SAST, DAST, IAST, RASP)}
% ─────────────────────────────────────────────────────────────────────────────

La sécurité applicative mobilise plusieurs familles de techniques de test,
chacune offrant un angle d'observation distinct sur la posture de sécurité
d'un système. Nous présentons ici les quatre approches majeures.

\subsection{SAST --- Static Application Security Testing}

\begin{definitionbox}{Définition --- SAST}
  Le \textbf{SAST} (test de sécurité statique) analyse le code source, le bytecode
  ou le code binaire d'une application \textit{sans l'exécuter}, afin d'identifier
  des vulnérabilités telles que les injections SQL, les dépassements de tampon,
  ou les mauvaises pratiques cryptographiques.
\end{definitionbox}

\paragraph{Principes de fonctionnement.}
Le SAST opère par analyse syntaxique et sémantique du code. Il construit
généralement un \textit{Abstract Syntax Tree} (AST) et effectue une analyse
du flux de données (\textit{taint analysis}) pour tracer le cheminement des
données non fiables depuis leurs sources jusqu'aux points d'utilisation
potentiellement dangereux (\textit{sinks}).

\paragraph{Avantages.}
\begin{itemize}
  \item Détection précoce des vulnérabilités, dès la phase de codage.
  \item Couverture exhaustive du code, y compris les chemins d'exécution rarement
    empruntés.
  \item Intégration aisée dans les IDE et les pipelines CI.
\end{itemize}

\paragraph{Limites.}
\begin{itemize}
  \item Taux élevé de faux positifs, nécessitant une validation humaine.
  \item Incapacité à détecter les vulnérabilités de configuration ou
    les failles logiques dépendant du contexte d'exécution.
  \item Difficulté à analyser les dépendances tierces compilées.
\end{itemize}

\paragraph{Outils représentatifs.}
SonarQube, Semgrep, Checkmarx, Fortify, Bandit (Python), SpotBugs (Java).

\subsection{DAST --- Dynamic Application Security Testing}

\begin{definitionbox}{Définition --- DAST}
  Le \textbf{DAST} (test de sécurité dynamique) évalue la sécurité d'une
  application \textit{en cours d'exécution}, en simulant des attaques réelles
  depuis l'extérieur, sans accès au code source. Il adopte la perspective
  d'un attaquant (\textit{black-box testing}).
\end{definitionbox}

\paragraph{Principes de fonctionnement.}
Le DAST interagit avec l'application via ses interfaces exposées (API REST,
interfaces web, gRPC, etc.), en injectant des charges malveillantes
(\textit{payloads}) et en analysant les réponses pour détecter des comportements
anormaux révélateurs de vulnérabilités : injections, authentification défaillante,
exposition de données sensibles, etc.

\paragraph{Avantages.}
\begin{itemize}
  \item Indépendant du langage et du framework applicatif.
  \item Détecte des vulnérabilités réelles dans l'environnement d'exécution,
    incluant celles issues de la configuration et des interactions entre composants.
  \item Faible taux de faux positifs par rapport au SAST.
\end{itemize}

\paragraph{Limites.}
\begin{itemize}
  \item Couverture partielle : ne peut tester que les fonctionnalités accessibles
    et les surfaces d'attaque exposées.
  \item Détection tardive dans le cycle de développement.
  \item Complexité accrue dans les architectures microservices, où la multitude
    des points d'entrée démultiplie la surface à couvrir.
\end{itemize}

\paragraph{Outils représentatifs.}
OWASP ZAP, Burp Suite, Nikto, Nuclei, Nessus.

\subsection{IAST --- Interactive Application Security Testing}

\begin{definitionbox}{Définition --- IAST}
  L'\textbf{IAST} (test de sécurité interactif) combine les approches statique
  et dynamique en instrumentant l'application à l'aide d'agents déployés dans
  son environnement d'exécution. Il analyse le comportement interne du code
  \textit{pendant} les tests fonctionnels.
\end{definitionbox}

\paragraph{Principes de fonctionnement.}
Des agents IAST sont intégrés directement dans le runtime de l'application
(JVM, CLR, interpréteur Python, etc.). Ils interceptent les flux d'exécution,
tracent les données et détectent les vulnérabilités en observant les interactions
réelles entre les composants, sans rejeter les requêtes légitimes.

\paragraph{Avantages.}
\begin{itemize}
  \item Très faible taux de faux positifs grâce à la corrélation
    contexte d'exécution / code source.
  \item Fonctionnement en temps réel, sans phase de scan dédiée.
  \item Adapté aux environnements de test d'intégration continue.
\end{itemize}

\paragraph{Limites.}
\begin{itemize}
  \item Surcharge de performance liée à l'instrumentation.
  \item Support limité à certains langages et frameworks.
  \item Couverture dépendante de l'exhaustivité des tests fonctionnels exécutés.
\end{itemize}

\paragraph{Outils représentatifs.}
Contrast Security, Seeker (Synopsys), Hdiv Security.

\subsection{RASP --- Runtime Application Self-Protection}

\begin{definitionbox}{Définition --- RASP}
  Le \textbf{RASP} (auto-protection applicative à l'exécution) est une technologie
  de sécurité intégrée directement dans une application ou son environnement
  d'exécution, capable de détecter et de bloquer des attaques en temps réel,
  sans intervention externe.
\end{definitionbox}

\paragraph{Principes de fonctionnement.}
Contrairement au DAST qui teste de l'extérieur, le RASP s'exécute à l'intérieur
de l'application. Il intercepte les appels système, les requêtes de base de données
et autres opérations sensibles, et peut les bloquer instantanément si un comportement
malveillant est détecté, tout en émettant des alertes vers le SIEM.

\paragraph{Avantages.}
\begin{itemize}
  \item Protection en production, y compris contre des attaques zero-day.
  \item Réduction de la dépendance aux mises à jour du WAF.
  \item Contexte d'exécution précis, minimisant les faux positifs.
\end{itemize}

\paragraph{Limites.}
\begin{itemize}
  \item Impact sur les performances applicatives.
  \item Risque de faux positifs entraînant des interruptions de service.
  \item Complexité de déploiement et de maintenance.
\end{itemize}

\paragraph{Outils représentatifs.}
Sqreen (Datadog), Imperva RASP, OpenRASP.

\subsection{Tableau comparatif des quatre approches}

\begin{table}[H]
  \centering
  \caption{Comparaison des approches de tests de sécurité}
  \label{tab:comparaison_tests}
  \renewcommand{\arraystretch}{1.4}
  \begin{tabularx}{\textwidth}{
    >{\bfseries}L{2cm} C{2.2cm} C{2.2cm} C{2.2cm} C{2.2cm}
  }
    \toprule
    \rowcolor{gray!15}
    Critère & SAST & DAST & IAST & RASP \\
    \midrule
    Phase & Codage / build & Test / staging & Test / intég. & Production \\
    \rowcolor{gray!10}
    Accès code & Requis & Non requis & Optionnel & Non requis \\
    Perspective & \textit{White-box} & \textit{Black-box} & \textit{Gray-box} & Interne \\
    \rowcolor{gray!10}
    Faux positifs & Élevés & Faibles & Très faibles & Faibles \\
    Couverture & Code complet & Surface exposée & Code exercé & Runtime \\
    \rowcolor{gray!10}
    Impact perf. & Aucun & Faible & Moyen & Moyen--élevé \\
    Microservices & Partiel & Complexe & Adapté & Adapté \\
    \bottomrule
  \end{tabularx}
\end{table}

% ─────────────────────────────────────────────────────────────────────────────
\section{Focus DAST : enjeux dans les architectures microservices}
% ─────────────────────────────────────────────────────────────────────────────

\subsection{Spécificités des architectures microservices}

Une architecture microservices est un style architectural dans lequel une
application est décomposée en un ensemble de services petits, autonomes et
faiblement couplés, chacun responsable d'une capacité métier précise.
Ces services communiquent entre eux via des API bien définies, généralement
sur des protocoles HTTP/REST ou gRPC, et peuvent être déployés indépendamment.

Cette architecture présente plusieurs caractéristiques qui complexifient
significativement l'application du DAST :

\subsubsection{Explosion de la surface d'attaque}

Dans une application monolithique, la surface d'attaque exposée est concentrée
sur un nombre limité de points d'entrée. Dans une architecture microservices,
chaque service expose potentiellement sa propre API, multipliant d'autant les
surfaces à tester. Un système composé de 50 microservices peut exposer plusieurs
centaines d'endpoints distincts, avec des contrats d'API hétérogènes.

\subsubsection{Dynamicité de l'infrastructure}

Les microservices s'exécutent dans des environnements hautement dynamiques
(Kubernetes, par exemple), où les instances de services apparaissent et
disparaissent selon la charge. Les adresses IP et les ports sont attribués
dynamiquement, rendant obsolète toute cible DAST définie statiquement.
Un scanner DAST doit s'adapter en temps réel à cette topologie mouvante.

\subsubsection{Communications interservices non exposées}

Une partie significative des communications dans une architecture microservices
est \textit{interne} : les échanges entre services du backend ne sont pas
directement accessibles depuis l'extérieur du cluster. Les outils DAST
conventionnels, conçus pour tester des interfaces web publiques, ne peuvent
pas atteindre ces flux de communication internes sans mécanismes spécifiques.

\subsubsection{Hétérogénéité des protocoles et des données}

Les microservices peuvent utiliser des protocoles variés : REST/HTTP, gRPC,
GraphQL, WebSocket, messaging asynchrone (Kafka, RabbitMQ). Les outils DAST
doivent être capables de générer des payloads adaptés à chaque protocole
et d'interpréter correctement les réponses correspondantes.

\subsection{Stratégies DAST adaptées aux microservices}

Face à ces défis, plusieurs stratégies permettent d'adapter le DAST aux
architectures microservices :

\paragraph{Intégration dans la pipeline CI/CD.}
Plutôt que d'effectuer des scans ponctuels, le DAST est intégré comme une
étape automatique de la pipeline, déclenchée à chaque déploiement en
environnement de staging. Des outils comme OWASP ZAP disposent de modes
``headless'' permettant cette intégration.

\paragraph{DAST orienté API.}
Des outils spécialisés dans le test de sécurité des API (comme 42Crunch,
Traceable AI, ou les extensions REST/gRPC de ZAP) permettent d'importer
les spécifications OpenAPI/Swagger et de générer automatiquement des cas
de test de sécurité couvrant l'ensemble des endpoints documentés.

\paragraph{Service Mesh comme point d'observation.}
L'adoption d'un Service Mesh ouvre une perspective nouvelle pour le DAST
dans les microservices. En opérant au niveau du sidecar proxy, le Service Mesh
fournit une visibilité totale sur les communications interservices, y compris
les flux internes au cluster, sans modification du code applicatif. Cette
caractéristique permet d'envisager un DAST \textit{interne}, capable de
tester la sécurité des échanges entre services --- une dimension absente
des approches DAST traditionnelles.

\begin{importantbox}
  \textbf{Synergie clé :} Le Service Mesh ne se substitue pas au DAST, mais
  il en élargit considérablement le périmètre d'action dans les architectures
  microservices. En fournissant un point d'interception centralisé et une
  observabilité totale du trafic interservices, il crée les conditions d'un
  DAST interne couvrant les vulnérabilités invisibles depuis l'extérieur
  du cluster.
\end{importantbox}

\subsection{Limites actuelles et pistes d'amélioration}

Malgré les évolutions récentes, plusieurs limites persistent dans l'application
du DAST aux architectures microservices :

\begin{itemize}
  \item \textbf{Couverture incomplète des flux asynchrones :} les communications
    via des bus de messages (Kafka, RabbitMQ) restent difficiles à intercepter
    et à tester dynamiquement.
  \item \textbf{Manque de standardisation :} l'absence d'un format universel de
    description des API interservices limite la capacité des outils DAST à
    générer automatiquement des cas de test pertinents.
  \item \textbf{Complexité de l'environnement de test :} reproduire fidèlement
    un environnement microservices complet à des fins de test reste coûteux
    en ressources et en configuration.
\end{itemize}

Ces limites motivent l'exploration d'une approche hybride intégrant le Service Mesh
comme infrastructure de sécurité complémentaire, objet des chapitres suivants.

% ─────────────────────────────────────────────────────────────────────────────
\section{Conclusion}
% ─────────────────────────────────────────────────────────────────────────────

Ce chapitre a posé les fondements conceptuels sur lesquels repose ce projet.
Nous avons d'abord introduit le paradigme DevSecOps, en retraçant son émergence
depuis le mouvement DevOps et en détaillant ses principes fondateurs, notamment
le \textit{Shift Left Security}. Nous avons ensuite dressé un panorama structuré
des quatre grandes familles de tests de sécurité --- SAST, DAST, IAST et RASP ---
en mettant en évidence leurs complémentarités et leurs limites respectives.

La section finale a approfondi les enjeux spécifiques du DAST dans les
architectures microservices, en identifiant les défis majeurs que soulève
cette combinaison : explosion de la surface d'attaque, dynamicité de
l'infrastructure, inaccessibilité des flux internes et hétérogénéité protocolaire.

\begin{tcolorbox}[
  colback=gray!8, colframe=darkgray,
  leftrule=4pt, toprule=0.4pt, bottomrule=0.4pt, rightrule=0.4pt,
  arc=2pt
]
  \textbf{Transition vers le Chapitre 2 :} Ces constats mettent en lumière la
  nécessité d'une infrastructure dédiée à la sécurisation des communications
  interservices. Le Service Mesh, présenté dans le chapitre suivant, constitue
  précisément une réponse à cette nécessité, en offrant une couche d'interception
  et de contrôle transparente sur l'ensemble du trafic du cluster.
\end{tcolorbox}
